\documentclass{article}
\usepackage{amsmath}

\begin{document}
    Buongiorno a tutti, oggi vi presento il mio lavoro di tesi sulla risommazione di grandi Logaritmi nella distribuzione
    del Thrust nell'annichilazione elettrone positrone. Vi presentero' il problema, risultati presenti in letterature e la 
    nostra continuazione del lavoro a ordini superiori in teoria delle perturbazioni.
    L'annichilazione elettrone positrone al livello zero in teoria delle perturbazioni e' questo diagramma di Feynman.
    Elettrone e positrone interagiscono tramite l'interazione elettromagnetica, mediata da un fotone virtuale che 
    successivamente decade sempre per interazione elettromagnetica in un quark e un antiquark. Questa coppia
    quark-antiquark dotati di carica di colore interagiscono forte ed emettono gluoni.
    
    Processi come questi sono ampiamente studiati negli acceleratori di particelle, che sono gradi "femto"-scopi che 
    indagare la natura piu' intima della materia, per e+e- 
    ad esempio sono stati studiato al LEP, operativo dal 1989 al 2000. Negli acceleratori vengono 
    misurate le proprieta' degli stati finali e vengono confrontate con le predizioni teoriche per
    testare il modello standard della fisica delle particelle, che e' la teoria delle particelle elementari di maggior successo 
    finora, ma anche per cercare nuova fisica oltre il modello standard e rispondere a domande
    ancora irrisolte come l'asimmetria materia-antimateria, la natura della materia oscura non 
    descritta dal modello standard.
    
    Ma prima di poter rispondere a queste domande, dobbiamo conoscere bene la teoria che abbiamo
    gia' a disposizione. Bisogna conoscere bene i parametri della teoria.
    
    Oggi vi presento uno dei modi con cui si puo' estrarre il parametro fondamentale della cromodinamica quantistica
    che e' la costante di accoppiamento forte, analogo alla costante di struttura fine dell'elettromagnetismo,
    il modo per farlo e' confrontare previsione per la distribuzione del thrust che dipende 
    dalla costante di accoppiamento forte con le misure sperimentali.
    
    La cromodinamica quantistica e' la teoria che descrive l'interazione forte, l'interazione che tiene confinato i quark dentro i protoni e neutroni.
    Gode di una proprieta' molto interessante che e' la liberta' asintotica, cioe' la costante di accoppiamento forte diventa piccola a energie elevate.
    Questo significa che possiamo fare calcoli perturbativi e ottenere risultati affidabilie confrontabili con le misure sperimentali.

    Il Thrust ha la definizione seguente: "formula"

    Dato una configurazione di momenti finali, immaginiamo il caso non banale con 3 momenti, e' definito come il massimo della della somma delle proiezioni dei momenti finali lungo un'asse diviso 
    per la somma in modulo dei momenti. L'asse che massimizza questa formula e' detto asse del Thrust e il valore ottenuto da questa formula con l'asse e' thrust e' detto Il
    valore del Thrust.
    
    Il Thrust e' una misura di quanto sono collimate le particelle finali, se il Thrust e' 1, le particelle sono perfettamente collimate lungo l'asse del Thrust, per valori minori di 1 le particelle sono meno collimate, fino 
    ad un valore minimo di 0.5 per particelle completamente isotrope, sfericamente distribuiti.

    Pero' il minimo del Thrust dipende dal numero di partoni o particelle (teoria ed esperimento usano la stessa definizione) 
    nello stato finale. Per N=3 il minimo e' assunto da questa configurazione trigonale planare con minimo 2/3, per N=4 il e' assunto 
    da questa configurazione tetraedrica minimo 1/sqrt(3). Mentre per N=5 in poi in letteratura non ho 
    trovato una soluzione analitica, ma solo una soluzione numerica estratta dai calcoli della sezione d'urto
    a ordine fisso che illustrero' in seguito.  


\end{document}